\section{Introduction}

In a competitive economy, a transaction changes more than the holdings of its two participants. Cash paid for a spare property can fund a rival's next investment. A deed can become an intermediary's bargaining instrument. Keeping a distressed opponent alive can prevent another player from inheriting a valuable estate. An agent evaluating only its own immediate acquisition would miss these consequences. An agent that talks about them has demonstrated a different, narrower achievement than one that chooses the best response to them.

This paper studies that distinction in complete games of Monopoly played by LLM agents. The scientific question is whether the agents' bargaining addresses the future options of competitors, and what the surrounding records establish about the use of that attention. We use \emph{strategic attention} for an explicit, state-linked consideration of another player's future holdings, finance, actions, or estate. It is an observable property of recorded conduct and explanations. We do not use it as a latent cognitive score. \emph{Strategic effectiveness} would require a criterion showing that the consideration improves a decision or outcome. That criterion is absent for most trades in this corpus.

The distinction matters because conversational sophistication can be read too generously or too dismissively. A seller who discusses a buyer's construction budget may be pricing a competitive threat, applying an ordinary markup, or rationalizing a demand after selecting it. A buyer who rejects a completing deed may be preserving liquidity, avoiding a rival's stronger completion, or making an error. Prices, stated reasons, and later results answer different parts of this problem. The purpose of following a complete game is to identify the stakes and the realized sequence without converting every winning trajectory into a good strategy.

Our central finding is that competitive consequences enter bargaining at several distances from the immediate exchange. Agents negotiate over a counterparty's building funds, a third party's resale opportunity, and a prospective creditor's acquisition. Some refuse their own group completion even after a proposer adjusts terms to address a cash objection. This is more specific than observing frequent trading or fluent strategic vocabulary. The same records also resist a stronger claim about coherent strategic valuation. Most counters to the examined deed exchanges reprice an unchanged bundle whether or not it completes a group. Aggregate higher asks for completing deeds become mixed in the small set of same-deed comparisons. In one focal negotiation, the correct concern about a rival's remaining cash accompanies an incorrect house-count calculation.

The twenty games contain 6,989 decisions and 526 negotiation threads between four fixed configurations. We examine the entire response population, then use eight contextual incidents to show what numerical categories omit. The incidents include a refused final-price claim and a settled one, refusal of one's own completion, correction followed by acceptance, a construction-financing denial appeal, selective disclosure during distress, withdrawal of a public assurance, and a payable debt resolved through a contradictory action. They were chosen retrospectively to expose different relationships between a statement and its consequence. They are neither a balanced model sample nor a conduct-frequency estimate.

The economic histories supply substantive checks on the interpretation. We add movement exposure and ownership timing to all sixty-three trade-completed groups. We reconstruct every episode reaching a debt menu, including 128 observed payments rather than studying only eliminations. We examine the positions immediately before the first estate, when the creditor already leads or ties in paid rent in fifteen games. These analyses make construction finance and creditor identity concrete objects of bargaining while constraining an explanation that assigns victory to inheritance alone.

The paper makes three contributions. First, it documents competitive-option bargaining in complete LLM game trajectories, with exact statements, current economic conditions, recipient responses, and later outcomes. Second, it tests simple aggregate interpretations against available within-record comparisons, including jointly stratified responses, same-deed asking prices, movement exposure, successful debt payments, and pre-estate positions. Third, it separates strategic attention, executable action, and realized consequence in a single observational account. That separation is useful because these configured participants do not operate under interchangeable response conditions.

This is a bounded study of existing play. Five seeds, cyclic rotations, one roster, uneven information, and response-specific recovery prevent causal attribution to a model or tactic. The broader research agenda concerns strategic reasoning, opponent modeling, negotiation and conduct, planning and memory, adaptation, economics and risk, collective dynamics, and validity. It also asks about strategic talk versus action, general versus multi-agent capability, and skill versus chance or weak opponents. Appendix~\ref{app:questions} maps all eleven questions to their measured fragments. Unmeasured capabilities remain unanswered questions.
